% Part: lambda-calculus
% Chapter: introduction
% Section: overview

\documentclass[../../../include/open-logic-section]{subfiles}

\begin{document}

\olfileid{lam}{int}{ovr}
\olsection{概述}

λ演算最初由 Alonzo Church 在20世纪30年代初设计，旨在作为构造性逻辑的基础，\emph{而非}作为可计算函数的模型。但人们很快发现它与其他可计算性定义（如图灵可计算函数和部分递归函数）是等价的。这个最初令人略感惊讶的事实，使得这种刻画更加有趣。

λ记法是一种便捷的方式，能直接通过定义它的符号表达式来指称该函数，而无需为它命名。我们不必说“设$f$是由$f(x) = x + 3$定义的函数”，而可以说“设$f$是函数$\lambd[x][(x + 3)]$”。换句话说，$\lambd[x][(x+3)]$仅仅是给自变元加3的那个函数的\emph{名字}。在这个表达式中，$x$是哑变元，或者说占位符：同一个函数同样可以表示为$\lambd[y][(y + 3)]$。即使周围有其他参数，这种记法依然适用。例如，假设$g(x, y)$是一个二元函数，而$k$是一个自然数。那么$\lambd[x][g(x,k)]$就是将任意~$x$映射到~$g(x, k)$的函数。

这种从符号表达式定义函数的方式称为\emph{λ抽象}。λ抽象的另一面是\emph{应用}：假设有一个函数$f$（例如定义在自然数上），我们可以将它\emph{应用}到任意值，比如2。当然，在传统记法中，我们把结果写作~$f(2)$。

将λ抽象与应用结合会产生什么结果？由此得到的表达式可以通过把被应用项“代入”所抽象的变元来进行简化。例如，
\[
(\lambd[x][(x + 3)])(2)
\]
可简化为~$2 + 3$。

到目前为止，我们所做的不过是为传统概念引入了新的记号。然而，λ演算代表了对集合论观点更为彻底的背离。在这个框架下：
\begin{enumerate}
\item 一切皆指称函数。
\item 函数可以用λ抽象来定义。
\item 任何事物都可以应用于任何事物。
\end{enumerate}

例如，如果$F$是λ演算中的一个项，那么$F(F)$总是假定是有意义的。这个自由的框架被称为\emph{无类型}λ演算，其中“无类型”意味着“对什么能应用于什么没有限制”。

\begin{digress}
还有一种\emph{有类型}λ演算，它是无类型版本的一个重要变体。尽管在许多方面有类型λ演算与无类型版本相似，但它更容易与经典的集合论框架协调，并且具有一些非常不同的性质。

λ演算的研究已被证明是理论计算机科学和编程语言设计中的核心部分。20世纪50年代由 John McCarthy 设计的 LISP，就是早期受这些思想影响的编程语言的一个例子。
\end{digress}

\end{document}